% ============================================================
%  PART IV  --  My Data
% ============================================================
\part{My Data}

% ── Chapter 12: My Data Overview ──────────────────────────────
\chapter{Mes données --- Aperçu général}
\label{ch:mydata}

Mes données sont le gestionnaire de ressources intégré de FairWindSK.
Il fournit une interface dédiée et optimisée pour toutes les données de navigation stockées sur le Signal~K
serveur : points de repère, itinéraires, régions, notes, cartes, historique des pistes et fichiers.

Open My Data by tapping the \key{My Data} icon (≡) in the Bottom Bar, or by tapping its tile on the
lanceur.

\screenshot{mydata-main-tabs}{Mes données --- barre des onglets montrant les sept onglets}{}

\begin{notebox}
Toutes les ressources de My Data sont stockées sur le serveur Signal~K, pas sur l'appareil FairWindSK.
Ils sont accessibles depuis n'importe quel périphérique connecté au même serveur Signal~K simultanément.
\end{notebox}

\section{Aperçu de l'onglet Mes données}

Mes données sont organisées en sept onglets.
Chaque onglet charge paresseusement --- il n'est créé que la première fois qu'il est tapé, gardant le démarrage rapide.

\rowcolors{2}{LtGray}{white}
\begin{longtable}{C{0.7cm} L{2.8cm} L{10.2cm}}
\toprule
\rowcolor{Navy}
{\sffamily\bfseries\color{white} Tab} &
{\sffamily\bfseries\color{white} Name} &
{\sffamily\bfseries\color{white} Contents} \\
\midrule
0 & Waypoints & Named positions: anchorages, hazards, marinas, fishing marks, race marks. \\
1 & Routes & Ordered sequences of waypoints forming planned passages. \\
2 & Regions & Geographic boundary areas --- SAR search regions, exclusion zones, anchorage limits. \\
3 & Notes & Free-text annotations attached to positions. \\
4 & Charts & Chart tile source catalogue for installed chart providers. \\
5 & Tracks & Historical GPS track --- the vessel's recorded position history. \\
6 & Files & Direct filesystem browser for files accessible through Signal~K. \\
\bottomrule
\end{longtable}

\section{La barre d'outils commune}

Les onglets 0--4 (Waypoints, Routes, Régions, Notes, Graphiques) partagent une interface commune :

\begin{itemize}
  \item A \textbf{title label} showing the resource type.
  \item A \textbf{search field} for real-time filtering of the list. Searching is case-insensitive and checks all visible columns simultaneously.
  \item A \textbf{scrollable table} showing all resources from the Signal~K server. Rows are 64\,px tall for comfortable touch selection. Alternating row colours aid readability. Column headers are clickable for sorting.
  \item A \textbf{toolbar} of icon buttons:
\end{itemize}

\rowcolors{2}{LtGray}{white}
\begin{longtable}{L{3cm} L{10.7cm}}
\toprule
\rowcolor{Navy}
{\sffamily\bfseries\color{white} Button} &
{\sffamily\bfseries\color{white} Action} \\
\midrule
Open ($\rightarrow$) & Opens the selected resource --- shows its details, opens an embedded map preview (Waypoints), or views the attachment (Notes). Accent-coloured; primary action. \\
Edit & Opens the resource editor drawer for the selected resource. \\
New (+) & Opens the resource editor drawer to create a new resource. \\
Delete & Deletes the selected resource after confirmation. Cannot be undone. \\
Import & Imports resources from a GeoJSON or JSON file. File picker opens in the Bottom Bar drawer. \\
Export & Exports the selected resource (or all resources if none selected) to a GeoJSON file. \\
Refresh & Forces a reload of all resources from the Signal~K server. \\
\bottomrule
\end{longtable}

\begin{notebox}
Sur Raspberry Pi OS (ARM Linux), la barre d'outils utilise des boutons texte uniquement pour éviter le rendu des icônes
des questions sur cette plate-forme.
La disposition et la fonctionnalité du bouton sont identiques; seule la présentation visuelle diffère.
\end{notebox}

% ── Chapter 13: Waypoints ─────────────────────────────────────
\chapter{Points de repère}
\label{ch:waypoints}

Les points de cheminement sont des positions nommées stockées sur le serveur Signal~K.
Ils sont utilisés pour les ancrages, les dangers, les marinas, les arrêts de carburant, les marques de pêche, les marques de course, et tout autre
autre position vers laquelle vous voulez enregistrer et naviguer.

\screenshot{mydata-waypoints-list}{Mes données --- Onglet Waypoints, liste avec nom, description, type, lat/lon, timestamp}{}

\section{Tableau des points de repère Colonnes}

Le tableau des points de repère montre six colonnes :

\rowcolors{2}{LtGray}{white}
\begin{longtable}{L{2.4cm} L{11.3cm}}
\toprule
\rowcolor{Navy}
{\sffamily\bfseries\color{white} Column} &
{\sffamily\bfseries\color{white} Description} \\
\midrule
Name & Display name of the waypoint, taken from the \code{name} field or from \code{feature.properties.name} in the GeoJSON resource. \\
Description & Free-text description of the waypoint. \\
Type & The waypoint type tag, e.g.\ \code{alarm-mob} (POB waypoints created by FairWindSK), \code{anchor}, \code{marina}, or any custom type. POB waypoints are visually distinguishable by their type. \\
Latitude & Latitude in your configured coordinate format (right-aligned). \\
Longitude & Longitude in your configured coordinate format (right-aligned). \\
Timestamp & ISO 8601 creation or last-update timestamp. \\
\bottomrule
\end{longtable}

\section{Recherche de Waypoints}

Le champ de recherche au-dessus du tableau filtre les points par n'importe quelle colonne visible.
Tapez une partie d'une valeur nom, description, type ou coordonnée.
La table se met à jour en temps réel lorsque vous tapez.

\section{Création d'un point de repère}

Tap \key{New} to open the Waypoint Editor drawer.


\subsection{Champs de l'éditeur}

\rowcolors{2}{LtGray}{white}
\begin{longtable}{L{3cm} L{10.7cm}}
\toprule
\rowcolor{Navy}
{\sffamily\bfseries\color{white} Field} &
{\sffamily\bfseries\color{white} Description} \\
\midrule
Name & The display name shown in the table, Top Bar waypoint widget, and chart plotter. \\
Description & Free-text notes about this position. \\
Type & Optional type tag. Leave blank for a generic waypoint. Use \code{anchor} for anchoring positions, \code{hazard} for obstructions, etc. \\
Latitude & Degrees with decimal minutes or decimal degrees, depending on the active coordinate format. Range: $-90$ to $+90$. \\
Longitude & Range: $-180$ to $+180$. \\
Altitude & Altitude in metres (optional). Leave as 0 if not relevant. \\
\bottomrule
\end{longtable}

After filling in the fields, tap \key{OK} in the editor.
Le point de cheminement est créé immédiatement sur le serveur Signal~K et apparaît dans la table.

\section{Modifier un Waypoint}

Select a waypoint in the table, then tap \key{Edit}.
Le même tiroir d'éditeur s'ouvre avec les valeurs existantes préremplies.
Modify any field and tap \key{OK}.
Le changement est écrit immédiatement à Signal~K.

\section{Supprimer un point de passage}

Select a waypoint and tap \key{Delete}.
Un tiroir de confirmation apparaît.
Tap \key{Delete} to confirm.
Le point de cheminement est retiré en permanence du serveur Signal~K.

\begin{warningbox}
La suppression ne peut être annulée.
POB waypoints (type \code{alarm-mob}) should generally be retained for accident reporting and debrief,
Non supprimé.
\end{warningbox}

\section{Ouverture d'un Waypoint --- la vue détaillée}

Select a waypoint and tap \key{Open} (or double-tap the row).
La vue détaillée de Waypoint s'ouvre.


La vue détaillée montre:

\begin{itemize}
  \item Full name, description, type, and formatted coordinates.
  \item An embedded Freeboard-SK map preview centred on the waypoint's position, if Freeboard-SK is installed on the connected Signal~K server. The map uses a JavaScript function to compute the view bounding box from the waypoint coordinates and zoom level.
  \item A \key{Navigate} button to set the waypoint as the active navigation destination.
\end{itemize}

\begin{notebox}
L'aperçu Freeboard-SK intégré est désactivé sur Raspberry Pi OS (ARM Linux) construit pour éviter
problèmes de performance avec des instances WebEngine simultanées sur cette plateforme.
Les coordonnées et les descriptions sont toujours affichées sous forme de texte, peu importe.
\end{notebox}

\section{Naviguer jusqu'à un point de passage}

From the detail view, tap \key{Navigate}.
FairWindSK calls \path{navigateToWaypoint} on the Signal~K client, which writes the waypoint's
resource path (e.g.\ \path{/resources/waypoints/\{uuid\}}) as the active navigation destination on
Signal ~K.
Immédiatement :

\begin{itemize}
  \item The Top Bar BTW, DTW, TTG, and ETA widgets begin showing values relative to this waypoint.
  \item Freeboard-SK and any other navigation application subscribed to \skpath{navigation.course} show the destination.
  \item The POB bearing and distance displays (when the POB bar is active) use this same destination.
\end{itemize}

\section{Importation de points de repère}

Tap \key{Import}. A file picker opens in the Bottom Bar drawer.
Select a GeoJSON file containing \code{Feature} objects with \code{Point} geometry.

- Oui.
\begin{enumerate}
  \item Reads the GeoJSON file.
  \item Extracts each \code{Point} feature as a waypoint.
  \item Creates or updates each waypoint on the Signal~K server (update if a matching UUID is found; create if not).
  \item Reports the count of imported waypoints in an informational drawer.
  \item Selects the first imported waypoint in the table.
\end{enumerate}

The import file filter accepts \code{.geojson} and \code{.json} files.

\section{Exportation de points de repère}

Tap \key{Export}. If a waypoint is selected, only that waypoint is exported.
Si rien n'est sélectionné, tous les points de repère sont exportés.
Un récupérateur de fichiers s'ouvre. Choisissez un nom de fichier et un emplacement.
The output is a valid GeoJSON \code{FeatureCollection} with \code{Point} geometry features.

% ── Chapter 14: Routes ────────────────────────────────────────
\chapter{Itinéraires}
\label{ch:routes}

Les itinéraires sont des séquences ordonnées de points formant des passages planifiés.
Ils sont généralement créés dans un graphique (Freeboard-SK) et stockés sur le serveur Signal~K.

\screenshot{mydata-routes-list}{Mes données --- Onglet Routes, liste avec nom, type de géométrie, nombre de points, horodatage}{}

\section{Chemins Tableau Colonnes}

\rowcolors{2}{LtGray}{white}
\begin{longtable}{L{2.4cm} L{11.3cm}}
\toprule
\rowcolor{Navy}
{\sffamily\bfseries\color{white} Column} &
{\sffamily\bfseries\color{white} Description} \\
\midrule
Name & Display name of the route. \\
Description & Free-text description or summary. \\
Geometry & The GeoJSON geometry type --- typically \code{LineString} for simple routes, \code{MultiLineString} for multi-leg routes with gaps, or \code{sar-spiral} / \code{sar-parallel} for SAR patterns created by the POB feature. \\
Points & The number of coordinate points in the route geometry. Displayed centred. \\
Timestamp & ISO 8601 creation or last-update timestamp. \\
\bottomrule
\end{longtable}

\section{Ouverture d'une route --- la vue détaillée}

Select a route and tap \key{Open}.


La vue détaillée montre:
\begin{itemize}
  \item Name, description, geometry type, and point count.
  \item A GeoJSON geometry preview if available.
  \item An \key{Activate} button to set the route's first waypoint as the active navigation destination.
\end{itemize}

\section{Création et édition des itinéraires}

\subsection{Créer une route}

Tap \key{New}. The Route Editor drawer opens.

\rowcolors{2}{LtGray}{white}
\begin{longtable}{L{3cm} L{10.7cm}}
\toprule
\rowcolor{Navy}
{\sffamily\bfseries\color{white} Field} &
{\sffamily\bfseries\color{white} Description} \\
\midrule
Name & Display name of the route. \\
Description & Free-text notes. \\
Coordinates (JSON) & The GeoJSON coordinate array for the route geometry. Enter as a JSON array of \code{[longitude, latitude]} pairs. \\
Geometry (JSON) & Advanced: edit the full GeoJSON geometry object if needed. \\
Properties (JSON) & Advanced: additional GeoJSON feature properties. \\
\bottomrule
\end{longtable}

\subsection{Activer une Route}

From the route detail view, tap \key{Activate}.
FairWindSK définit le premier point de départ de l'itinéraire comme étant la destination de navigation active Signal~K.
The Autopilot Route mode (Chapter~\ref{ch:autopilot}) will then follow the active route automatically.

\section{Routes d'importation et d'exportation}

Import and export use the same GeoJSON \code{FeatureCollection} format as Waypoints.
Le type de géométrie dans le GeoJSON détermine comment la route est stockée sur Signal~K.

\begin{tipbox}
SAR search routes created by the POB feature (Chapter~\ref{ch:pob}) appear in the Routes tab
automatiquement après la création.
They have geometry types \code{sar-spiral} and \code{sar-parallel} and can be exported as GeoJSON
pour les relevés de voyage.
\end{tipbox}

% ── Chapter 15: Regions ───────────────────────────────────────
\chapter{Régions}
\label{ch:regions}

Les régions sont des zones limites géographiques définies comme étant des polygones GeoJSON.
Ils sont utilisés pour les zones de recherche SAR (automatiquement créées par la fonction POB), zone marine protégée
limites, zones d'exclusion, limites d'ancrage et zones de pêche.

\screenshot{mydata-regions-list}{Mes données --- Onglet Régions, liste avec nom, type de géométrie, horodatage}{}

\section{Régions Tableau Colonnes}

\rowcolors{2}{LtGray}{white}
\begin{longtable}{L{2.4cm} L{11.3cm}}
\toprule
\rowcolor{Navy}
{\sffamily\bfseries\color{white} Column} &
{\sffamily\bfseries\color{white} Description} \\
\midrule
Name & Display name of the region. \\
Description & Free-text description. \\
Geometry & The GeoJSON geometry type --- typically \code{Polygon} or \code{MultiPolygon}. SAR regions use \code{Polygon}. \\
Timestamp & ISO 8601 creation or last-update timestamp. \\
\bottomrule
\end{longtable}

\section{Création d'une région}

Tap \key{New}. The Region Editor drawer opens.

\rowcolors{2}{LtGray}{white}
\begin{longtable}{L{3cm} L{10.7cm}}
\toprule
\rowcolor{Navy}
{\sffamily\bfseries\color{white} Field} &
{\sffamily\bfseries\color{white} Description} \\
\midrule
Name & Display name. \\
Description & Free-text notes. \\
Coordinates (JSON) & The GeoJSON coordinate array for the polygon. Polygons require an outer ring (first and last coordinate identical to close the ring). \\
Geometry (JSON) & Advanced: full GeoJSON geometry object. \\
Properties (JSON) & Advanced: additional GeoJSON feature properties. \\
\bottomrule
\end{longtable}

\begin{notebox}
Les zones de recherche SAR créées par la fonction POB apparaissent automatiquement dans l'onglet Régions.
Their \code{properties} include the \code{pobUuid} linking them to the originating POB incident.
\end{notebox}

% ── Chapter 16: Notes ─────────────────────────────────────────
\chapter{Annexe}
\label{ch:notes}

Les notes sont des annotations en texte libre attachées aux positions géographiques.
Ils fonctionnent comme des notes collantes sur le graphique --- utiles pour enregistrer les conseils de pilotage, les dangers locaux non
sur la carte, les détails de la marina ou les observations de l'équipage.

\screenshot{mydata-notes-list}{Mes données --- Onglet Notes, vue de liste avec titre, description, href, type MIME, timestamp}{}

\section{Notes Tableau Colonnes}

Les notes ont un ensemble de colonnes distinctes reflétant leur nature axée sur l'attachement :

\rowcolors{2}{LtGray}{white}
\begin{longtable}{L{2.4cm} L{11.3cm}}
\toprule
\rowcolor{Navy}
{\sffamily\bfseries\color{white} Column} &
{\sffamily\bfseries\color{white} Description} \\
\midrule
Title & The display name of the note (the first column is labelled \textit{Title} for notes, not \textit{Name}, to reflect the note-taking context). \\
Description & Summary text. For imported text files, this is set to the first 120 characters of the file content. \\
Href & An optional file path or URL linking the note to an attachment --- a PDF pilot book page, an image, a text file. When present, tapping \key{Open} opens the attachment directly in the built-in viewer. \\
MIME Type & The MIME type of the attachment, e.g.\ \code{application/pdf}, \code{image/jpeg}, \code{text/plain}. Used to route the file to the correct viewer. \\
Timestamp & ISO 8601 creation or last-update timestamp. \\
\bottomrule
\end{longtable}

\section{Création d'une note}

Tap \key{New}. The Note Editor drawer opens.


\rowcolors{2}{LtGray}{white}
\begin{longtable}{L{3cm} L{10.7cm}}
\toprule
\rowcolor{Navy}
{\sffamily\bfseries\color{white} Field} &
{\sffamily\bfseries\color{white} Description} \\
\midrule
Title & The note title (displayed in the \textit{Title} column). \\
Description & Main text content of the note. \\
Href & Optional path or URL to an attachment file. Enter an absolute path (e.g.\ \path{/home/pi/charts/harbour.pdf}) or a URL. \\
MIME Type & Optional MIME type for the attachment. FairWindSK uses this to select the correct viewer. If left blank, FairWindSK detects the type from the file extension. \\
Has position & Check this box to attach a geographic position to the note. \\
Latitude & Latitude when \textit{Has position} is checked. \\
Longitude & Longitude when \textit{Has position} is checked. \\
\bottomrule
\end{longtable}

\section{Ouverture d'une note}

\subsection{Notes sans pièce jointe}

Tapping \key{Open} opens the note editor in read-only mode, showing the title, description, and position.

\subsection{Notes avec une pièce jointe (ensemble de mailles)}

Tapping \key{Open} immediately opens the attachment:

\begin{itemize}
  \item \textbf{Local file that exists:} opens in the built-in FileViewer (PDF, image, or text viewer depending on the MIME type).
  \item \textbf{HTTP/HTTPS URL:} opens in an embedded web view inside the Bottom Bar horizontal drawer.
  \item \textbf{Local directory path:} shows a message directing you to use My Data~> Files to browse the folder (single-window model prevents opening a separate file manager).
\end{itemize}

\section{Importation de notes}

L'onglet Notes accepte un ensemble plus large de types de fichiers d'importation que les autres onglets.
En plus de GeoJSON/JSON:

\begin{itemize}
  \item \textbf{Text files} (\code{.txt}, \code{.md}, \code{.html}): FairWindSK creates a note with the filename as the title and the first 120 characters of the file content as the description, setting the \code{href} to the absolute file path and detecting the MIME type automatically.
  \item \textbf{Images} (\code{.jpg}, \code{.jpeg}, \code{.png}): Creates a note with the filename as title and sets the \code{href} and \code{mimeType} accordingly.
  \item \textbf{PDF files}: Creates a note pointing to the PDF file via \code{href}.
\end{itemize}

The import filter text reads: \textit{Notes and attachments (*.geojson *.json *.txt *.md *.markdown *.html *.htm *.pdf *.jpg *.jpeg *.png)}.

% ── Chapter 17: Charts ────────────────────────────────────────
\chapter{Graphiques}
\label{ch:charts}

L'onglet Graphiques énumère les sources disponibles de carreau de cartes enregistrées sur le serveur Signal~K.
Il gère le catalogue source des cartes --- pas l'affichage réel des cartes, qui se produit à l'intérieur
Freeboard-SK ou un autre traceur de cartes fonctionnant dans la zone d'application.

\screenshot{mydata-charts-list}{Mes données --- onglet Graphiques, vue de liste avec nom, format, identifiant, URL, horodatage}{}

\section{Graphiques Tableau Colonnes}

\rowcolors{2}{LtGray}{white}
\begin{longtable}{L{2.4cm} L{11.3cm}}
\toprule
\rowcolor{Navy}
{\sffamily\bfseries\color{white} Column} &
{\sffamily\bfseries\color{white} Description} \\
\midrule
Name & Display name of the chart source. \\
Description & Free-text description. \\
Format & The chart tile format. Supported formats include: \code{mbtiles} (MBTiles SQLite database), \code{kap} (BSB/KAP raster charts), \code{gemf} (GEMF tile archive), \code{tilemapresource} (TileMapResource XML), \code{package} (ZIP chart package). \\
Identifier & The chart identifier used by Signal~K and chart plotters to reference this source. \\
Chart URL & The local filesystem path or URL where the chart tiles are served. \\
Timestamp & ISO 8601 registration timestamp. \\
\bottomrule
\end{longtable}

\section{Ajout d'une source graphique}

Tap \key{New}. The Chart Editor drawer opens.

\rowcolors{2}{LtGray}{white}
\begin{longtable}{L{3cm} L{10.7cm}}
\toprule
\rowcolor{Navy}
{\sffamily\bfseries\color{white} Field} &
{\sffamily\bfseries\color{white} Description} \\
\midrule
Name & Display name for this chart source. \\
Description & Free-text description. \\
Identifier & Unique identifier string used by chart plotters to reference this source. \\
Chart format & The format of the chart tiles. FairWindSK detects the format from the file extension when importing. \\
Chart URL & Path or URL to the chart tiles. For MBTiles, this is the absolute path to the \code{.mbtiles} file. \\
Tilemap URL & For \code{tilemapresource} format only: the path to the \code{tilemapresource.xml} file. The Chart URL is then set to the directory containing it. \\
Chart region & Optional: a GeoJSON bounding box or geometry describing the geographic coverage of this chart. \\
Chart scale & Optional: the chart scale denominator (e.g.\ \code{50000} for 1:50,000). \\
Chart layers & Optional: comma-separated list of layer names for tile sources that support multiple layers. \\
Chart bounds & Optional: GeoJSON coordinates describing the chart bounds. \\
\bottomrule
\end{longtable}

\section{Importation d'une source graphique}

Tap \key{Import}. Select a chart file.
FairWindSK détecte le format graphique à partir de l'extension de fichier et construit un objet de ressource graphique
automatiquement.
For \code{tilemapresource.xml} files, FairWindSK sets both the tile URL (the directory) and the
URL carlemap (le chemin de fichier XML).

L'import crée ou met à jour l'enregistrement du graphique sur le serveur Signal~K.

\section{Exportation du catalogue des graphiques}

Tap \key{Export}. The chart catalogue is exported as a JSON manifest document with structure:

\begin{lstlisting}
{
  "type": "FairWindSKChartPackage",
  "charts": [ { ... }, { ... } ]
}
\end{lstlisting}

Ce manifeste peut être importé sur une autre installation FairWindSK pour reproduire le catalogue de cartes
la configuration (elle décrit les sources des cartes, et non les données des tuiles).

% ── Chapter 18: Tracks ────────────────────────────────────────
\chapter{Voies --- Historique}
\label{ch:tracks}

L'onglet Voies indique la trajectoire GPS historique du navire --- la séquence de position enregistrée
s'arrange avec le temps.
Les pistes sont générées automatiquement par Signal~K lorsque le plugin des pistes est activé.

\screenshot{mydata-tracks-list}{Mes données --- Tab des pistes, tableau des points de piste avec le temps, lat/lon, altitude}{}

\section{Colonnes de tableau de voie}

Le tableau de la voie montre les différentes corrections de position (points) par ordre chronologique:

\rowcolors{2}{LtGray}{white}
\begin{longtable}{L{2.4cm} L{11.3cm}}
\toprule
\rowcolor{Navy}
{\sffamily\bfseries\color{white} Column} &
{\sffamily\bfseries\color{white} Description} \\
\midrule
Time & The timestamp of the position fix, formatted in local time using the system locale. \\
Latitude & Latitude to 6 decimal places. \\
Longitude & Longitude to 6 decimal places. \\
Altitude & Altitude in metres (if available from the GPS fix; otherwise shown as 0). \\
\bottomrule
\end{longtable}

\section{Enregistrement en direct de la piste}

Si Signal~K fournit activement des mises à jour de position, de nouveaux points de piste sont ajoutés au tableau dans
en temps réel.
La table défile pour afficher automatiquement le point le plus récent.

\section{Affichage d'un point de piste}

Select a row and tap \key{Open}. A read-only dialog opens showing the timestamp, latitude, longitude,
et l'altitude pour ce point en détail.

\subsection{Modifier un point de piste}

Select a row and tap \key{Edit}. The Track Point Editor opens.

\rowcolors{2}{LtGray}{white}
\begin{longtable}{L{3cm} L{10.7cm}}
\toprule
\rowcolor{Navy}
{\sffamily\bfseries\color{white} Field} &
{\sffamily\bfseries\color{white} Description} \\
\midrule
Timestamp & Date and time of the fix. Uses a date/time picker with up/down arrows. \\
Latitude & Latitude (8 decimal places, range $-90$ to $+90$). \\
Longitude & Longitude (8 decimal places, range $-180$ to $+180$). \\
Altitude & Altitude in metres. \\
\bottomrule
\end{longtable}

\section{Ajouter un point de piste manuellement}

Tap \key{New}. The same editor opens with default values (current time and position 0,0).
Enter the correct values and tap \key{OK}.
Le point est inséré dans la piste dans l'ordre chronologique.

\section{Supprimer un point de piste}

Select a row and tap \key{Delete}. Confirm when prompted. The point is removed from the track.

\section{Exportation de la voie}

Tap \key{Export} to save the complete track as a GeoJSON \code{LineString} or
\code{FeatureCollection}:

\begin{lstlisting}
{
  "type": "FeatureCollection",
  "features": [
    {
      "type": "Feature",
      "geometry": {
        "type": "LineString",
        "coordinates": [[lon, lat, alt], ...]
      }
    }
  ]
}
\end{lstlisting}

The default export filename is \path{tracks-history.geojson}.

\begin{tipbox}
Exporter la piste après chaque passage et la stocker avec le carnet de voyage.
Le fichier GeoJSON peut être importé dans la plupart des outils SIG, traceurs de cartes et logiciel de planification de voyage
pour examen et analyse.
\end{tipbox}

\section{Importation d'une piste}

Tap \key{Import} and select a GeoJSON file containing a \code{LineString} geometry.
FairWindSK extrait chaque coordonnée dans la ligne comme un point de piste avec un horodatage synthétique,
les annexer à la piste actuelle dans l'ordre.

\section{Utilisations de l'onglet Track}

\begin{itemize}
  \item \textbf{Passage review:} compare the actual track against the planned route. Identify deviations, traffic avoidance manoeuvres, and current effects.
  \item \textbf{Anchorage assessment:} see the swing circle after anchoring overnight.
  \item \textbf{Safety investigation:} reconstruct the vessel's path before and during an incident for formal reporting.
  \item \textbf{Logbook cross-reference:} correlate positions at specific times with logbook entries.
\end{itemize}

% ── Chapter 19: Files ─────────────────────────────────────────
\chapter{Fichiers}
\label{ch:files}

L'onglet Fichiers est un navigateur direct du système de fichiers pour les fichiers et les répertoires accessibles sur l'appareil
Il dirige FairWindSK.
Il fournit des opérations de navigation, de recherche, de visualisation et de gestion de fichiers (copie, découpe, coller, renommer,
supprimer, nouveau dossier) sans quitter l'interface FairWindSK.

\screenshot{mydata-files-browser}{Mes données --- Onglet Fichiers, navigateur répertoire avec barre d'outils}{}

\section{Mise en page de l'interface}

L'onglet Fichiers est divisé en trois zones :

\begin{itemize}
  \item \textbf{Toolbar} at the top: navigation and action buttons, path bar, and search field.
  \item \textbf{File list} in the centre: the current directory contents as a tree view with columns for Name, Size, and Date. Row height is 58\,px for comfortable touch selection. Alternating row colours aid readability.
  \item \textbf{Actions row} below the toolbar: file operation buttons that appear when items are selected.
\end{itemize}

\section{Navigation}

\subsection{Navigation des répertoires}

\rowcolors{2}{LtGray}{white}
\begin{longtable}{L{3cm} L{10.7cm}}
\toprule
\rowcolor{Navy}
{\sffamily\bfseries\color{white} Action} &
{\sffamily\bfseries\color{white} How to trigger} \\
\midrule
Enter a folder & Double-tap a folder row in the file list, or select and tap \key{Open}. \\
Go up one level & Tap \key{↑} (Up) button in the toolbar. \\
Navigate by path & Tap the path bar at the top of the toolbar, type an absolute path, and press Return. \\
Go home & Tap \key{Home} button. Returns to the home directory of the running user. \\
\bottomrule
\end{longtable}

\subsection{Tri des colonnes}

Appuyez sur n'importe quel en-tête de colonne pour trier par cette colonne.
La direction de tri oscille entre la montée et la descente sur les robinets suivants.
Le tri par défaut est par Nom, ascendant.

\section{Recherche de fichiers}

\subsection{Recherche fondamentale}

Tap \key{Search} or press Return in the search field.
FairWindSK recherche récursivement l'arborescence du répertoire actuel pour les fichiers et sous-répertoires dont les noms
contient le texte de recherche.

Les résultats sont affichés dans une liste à plat dans la vue de fichier.
Appuyez sur un résultat pour le sélectionner; double-tap pour ouvrir.
La recherche continue jusqu'à la racine de l'arborescence du répertoire parcouru.


\subsection{Options de recherche}

Tap \key{Options} (or the filter button) to expand the search options:

\rowcolors{2}{LtGray}{white}
\begin{longtable}{L{3.5cm} L{10.2cm}}
\toprule
\rowcolor{Navy}
{\sffamily\bfseries\color{white} Option} &
{\sffamily\bfseries\color{white} Description} \\
\midrule
Case Sensitive & When enabled, the search matches the exact capitalisation of the search text. Default: off (case-insensitive). \\
Search Hidden & When enabled, hidden files and directories (names beginning with \code{.}) are included in search results. Default: off. \\
Search System & When enabled, system files and directories (e.g.\ \code{/proc}, \code{/sys} on Linux) are included. Default: off. \\
\bottomrule
\end{longtable}

The search options button label updates to show a compact summary of active options (e.g.\
\textit{Options: case, hidden}).

\section{Actions de fichiers}

\subsection{Ouverture des fichiers}

Select a file and tap \key{Open}, or double-tap the row.

FairWindSK détecte le type de fichier en utilisant la détection de type MIME et l'oriente vers le visionneur approprié:

\rowcolors{2}{LtGray}{white}
\begin{longtable}{L{4cm} L{9.7cm}}
\toprule
\rowcolor{Navy}
{\sffamily\bfseries\color{white} File type} &
{\sffamily\bfseries\color{white} Viewer} \\
\midrule
Images (\code{.jpg}, \code{.jpeg}, \code{.png}, \code{.gif}) & Built-in image viewer with pinch-to-zoom and pan. \\
PDF (\code{.pdf}) & Built-in PDF viewer, page by page. \\
Text files (\code{.txt}, \code{.md}, \code{.log}, \code{.html}) & Built-in text viewer, formatted as plain text or HTML. \\
Other file types & A message is shown; the file cannot be previewed but can be managed. \\
\bottomrule
\end{longtable}

Les répertoires s'ouvrent dans le navigateur de répertoire (navigant en eux) plutôt que dans un visionneur.
Essayer d'ouvrir un chemin de répertoire à partir d'un href Notes affiche un message vous ordonnant d'utiliser les Fichiers
onglet pour le parcourir à la place (contrainte modèle à guichet unique).


\subsection{Renaming des fichiers et dossiers}

Select an item and tap \key{Rename}.
Une boîte de dialogue s'ouvre dans le tiroir horizontal de la barre inférieure.
Enter the new name and tap \key{OK}.
Le fichier ou répertoire est rebaptisé immédiatement sur le système de fichiers.

\subsection{Création d'un nouveau dossier}

Tap \key{New Folder}.
Une boîte de dialogue d'entrée s'ouvre.
Enter the folder name and tap \key{OK}.
Le nouveau dossier est créé dans le répertoire actuel.

\subsection{Copier, couper et coller}

\begin{enumerate}
  \item Select one or more files or folders.
  \item Tap \key{Copy} to mark them for copying, or \key{Cut} to mark them for moving.
  \item Navigate to the destination directory.
  \item Tap \key{Paste}.
\end{enumerate}

FairWindSK effectue des contrôles de sécurité avant les opérations de pâte:

\begin{itemize}
  \item Copying a directory into itself or one of its subdirectories is blocked.
  \item Pasting into the source directory of a cut operation shows a warning.
  \item The self-paste safety check runs a self-test at startup to verify the detection logic is working correctly.
\end{itemize}

Les opérations de copie et de déplacement des répertoires sont récursives --- l'ensemble du sous-arbre est copié ou déplacé.

\subsection{Suppression des fichiers et dossiers}

Select one or more items and tap \key{Delete}.
Une boîte de dialogue de confirmation s'ouvre dans le tiroir horizontal de la barre inférieure.
Confirmer de supprimer.

\begin{warningbox}
La suppression de fichier est permanente.
Il n'y a pas de bac de recyclage ou d'annulation pour les opérations de fichiers.
Prenez particulièrement soin de supprimer les répertoires --- tous les contenus sont supprimés récursivement.
\end{warningbox}

\section{Colonnes de la liste de fichiers}

La liste de fichiers montre trois colonnes :

\rowcolors{2}{LtGray}{white}
\begin{longtable}{L{2cm} L{11.7cm}}
\toprule
\rowcolor{Navy}
{\sffamily\bfseries\color{white} Column} &
{\sffamily\bfseries\color{white} Description} \\
\midrule
Name & Filename or directory name, with an icon indicating the file type (folder, image, document, etc.). \\
Size & File size in human-readable units (B, KB, MB, GB). Directories show the size of their immediate contents. \\
Date & Last modification timestamp. The \textit{Type} column (column index 2) is hidden by default. \\
\bottomrule
\end{longtable}

\section{Utilisations pratiques de l'onglet Fichiers}

\begin{itemize}
  \item Browse pilot book excerpts stored as PDFs on the Raspberry Pi.
  \item View passage planning images (harbour photos, aerial views, approach charts).
  \item Access voyage log files or NMEA recordings.
  \item Manage chart tile archives before registering them in the Charts tab.
  \item Review diagnostic log files generated by FairWindSK (see Chapter~\ref{ch:settings-system}).
\end{itemize}
